\documentclass[11pt]{article}
\usepackage{amsmath,amssymb,amsthm}
\usepackage[margin=1in]{geometry}
\usepackage{hyperref}

\newtheorem{theorem}{Theorem}
\newtheorem{proposition}{Proposition}
\newtheorem{remark}{Remark}

\newcommand{\bbF}{\mathbb{F}}
\newcommand{\bbQ}{\mathbb{Q}}
\newcommand{\bbZ}{\mathbb{Z}}
\newcommand{\cO}{\mathcal{O}}
\newcommand{\rhob}{\bar{\rho}}
\newcommand{\Gal}{\mathrm{Gal}}
\newcommand{\Frob}{\mathrm{Frob}}
\newcommand{\Jac}{\mathrm{Jac}}
\newcommand{\Tr}{\mathrm{Tr}}
\newcommand{\Nm}{\mathrm{Nm}}
\newcommand{\End}{\mathrm{End}}
\newcommand{\chicyc}{\chi_{11}}
\DeclareMathOperator{\GSp}{GSp}

\title{A mod-11 modularity match over $\bbQ(\sqrt5)$ \\ from a $(1,11)$-polarized abelian surface}
\author{(working note)}
\date{June 2026}

\begin{document}
\maketitle

\begin{abstract}
We exhibit computational evidence for Serre's modularity conjecture over the real quadratic
field $K=\bbQ(\sqrt5)$. Starting from a point of the Gross--Popescu moduli of
$(1,11)$-polarized abelian surfaces, we construct a genus-2 curve $C$ over $K$ whose Jacobian
$A=\Jac(C)$ is geometrically simple ($\End_{\bar K}A=\bbZ$) with good reduction at both primes
above $11$. The mod-$11$ representation $A[11]$ splits over $K$ as a sum of two irreducible
$2$-dimensional pieces $\rho_1,\rho_2$ with cyclotomic determinant --- a splitting forced not
by extra endomorphisms but by the polarization, since $11$ is the polarization degree. Working
entirely mod $11$ inside a space of Hilbert modular forms computed via a definite quaternion
algebra, we exhibit a Hilbert modular eigenform of parallel weight $2$ and level
$\mathfrak p\cdot(2)$ over $K$ ($\mathfrak p$ the bad prime of $A$ of norm $66179$) realizing
one of the pieces up to the quadratic twist $\chi_{K(\sqrt d)}$, $d=66179(5-\sqrt5)/2$
(ramified at $5$ and $66179$), verified over $24$ primes --- concrete evidence for Serre
modularity over $K$. The modular form itself supplies the global labeling of the two pieces.
\end{abstract}

\section{Introduction}

Serre's modularity conjecture, a theorem over $\bbQ$ (Khare--Wintenberger), remains largely
open over totally real fields. In its Hilbert-modular form it predicts that every odd,
irreducible, continuous representation
$\rhob\colon \Gal(\bar K/K)\to \mathrm{GL}_2(\bar\bbF_\ell)$ of the absolute Galois group of a
totally real field $K$ arises from a Hilbert modular eigenform. This note records a concrete
instance over $K=\bbQ(\sqrt5)$ with $\ell=11$, obtained by the following pipeline:
\[
\text{$(1,11)$ Gross--Popescu point} \;\longrightarrow\; \text{abelian surface }A/K
\;\longrightarrow\; A[11]=\rho_1\oplus\rho_2 \;\longrightarrow\; \text{Hilbert modular form}.
\]
The point of the construction is that the abelian surface comes with a canonical mod-$11$
structure (the $(1,11)$-polarization), which is exactly what makes $A[11]$ reducible into
$2$-dimensional pieces and supplies interesting $\rhob$'s to test.

\section{The abelian surface}

Let $K=\bbQ(\sqrt5)$, $\cO_K=\bbZ[\tfrac{1+\sqrt5}2]$, of class number $1$ and discriminant
$5$. The prime $11$ splits in $K$.

The Gross--Popescu construction (\texttt{arXiv:math/9902017}) realizes a $(1,11)$-polarized
abelian surface in $\mathbb P^{10}$ from a point of the sextic ``odd $2$-torsion'' locus
$D_2\subset\mathbb P^4$. Inverting the associated theta map numerically yields a period matrix
$\tau$, hence (via its principally polarized partner) a genus-2 curve $C/K$. We work with the
curve {\bf Curve A}, with absolute Igusa--Siegel invariants (normalized $J_2=1$)
\[
\Big[\,1,\ \tfrac{243125\sqrt5-482787}{1459240},\
\tfrac{-54798934\sqrt5+127710391}{2229718720},\
\tfrac{182422196780\sqrt5-406668692089}{8517525510400},\
\tfrac{38134182372761\sqrt5-85270624049895}{65073894899456000}\,\Big].
\]

\begin{proposition}\label{prop:A}
$A=\Jac(C)$ satisfies:
\begin{enumerate}
\item[(i)] $A$ is geometrically simple with $\End_{\bar K}(A)=\bbZ$ (no real multiplication);
\item[(ii)] $A$ has good reduction at both primes of $K$ above $11$;
\item[(iii)] away from $11$, the conductor of $A$ has odd part $\mathfrak p$, a prime of $K$
of norm $66179$, at which $A$ has multiplicative reduction (conductor exponent $1$);
\item[(iv)] $A$ also has (model-)bad reduction at the inert prime $(2)$.
\end{enumerate}
\end{proposition}

These are verified by computation: (i) via the geometric endomorphism algebra (CHIMP), (ii)
and (iii) via the Igusa valuation criterion and conductor exponents.

\section{The mod-11 representation}

Let $\rhob\colon\Gal(\bar K/K)\to\GSp_4(\bbF_{11})$ be the representation on $A[11]$, and let
$\chicyc$ denote the mod-$11$ cyclotomic character. We summarize the structure, all of which
is verified over the $290$ good primes $P$ of $K$ with $N(P)\le 2000$ by reducing $C$ and
factoring the Frobenius characteristic polynomial $h_P(x)\in\bbF_{11}[x]$.

\begin{proposition}\label{prop:struct}
\begin{enumerate}
\item[(a)] For every such $P$, $h_P$ factors over $\bbF_{11}$ as a product of two quadratics:
its factor type is always $[1,1,2]$ or $[1,1,1,1]$, and never $[2,2]$, $[1,3]$, or $[4]$.
\item[(b)] $\rhob$ has no $1$-dimensional $K$-subquotient: there is no Galois character
$\lambda$ that is, at every prime, a root of the ``linear part'' of $h_P$ (an exhaustive search
over Hecke characters of $K$ of order dividing $10$ and every plausible conductor gives at best
$90/132$ agreement).
\item[(c)] Consequently $A[11]$ is, up to semisimplification, a sum
$\rho_1\oplus\rho_2$ of two $2$-dimensional representations, each with
$\det\rho_i=\chicyc$ (hence odd). The splitting is supplied by the polarization: the kernel
$K(\mathcal L)\cong(\bbZ/11)^2$ of the $(1,11)$-polarization is a $K$-rational $2$-dimensional
subspace of $A[11]$. The prime $11$ is special precisely because it is the polarization degree;
at other primes $\ell$ the image of $\Jac(C)[\ell]$ is the full $\GSp_4(\bbF_\ell)$.
\item[(d)] $\rho_1,\rho_2$ are irreducible over $K$, are not isomorphic, and are not quadratic
twists of one another; their image is large (not dihedral/CM: the surface trace vanishes at
only $\approx 10\%$ of primes).
\end{enumerate}
\end{proposition}

Because $\det\rho_i=\chicyc$, the two traces are the roots of an explicit quadratic. Writing
$h_P(x)=x^4+Ax^3+Bx^2+Cx+D$ (so $D=N(P)^2$ and $C=N(P)\,A$), one has
\begin{equation}\label{eq:fingerprint}
\{\,\mathrm{tr}\,\rho_1(\Frob_P),\ \mathrm{tr}\,\rho_2(\Frob_P)\,\}
=\text{roots of }\ y^2 + A\,y + \big(B-2N(P)\big)\ \ (\mathrm{mod}\ 11).
\end{equation}
We call $\big(\Tr_P,\Nm_P\big):=\big(-A,\ B-2N(P)\big)\bmod 11$ the {\it fingerprint} of
$\{\rho_1,\rho_2\}$ at $P$. A sample:
\[
\begin{array}{c|cccccccc}
N(P) & 49 & 169 & 29 & 29 & 31 & 31 & 41 & 41\\\hline
\Tr_P & 9 & 6 & 7 & 9 & 5 & 0 & 9 & 2\\
\Nm_P & 9 & 5 & 1 & 7 & 4 & 10 & 0 & 3
\end{array}
\]

\section{The modularity match}

By good reduction at $11$ (Proposition~\ref{prop:A}(ii)) each $\rho_i$ is finite-flat at $11$,
so the predicted Serre weight is parallel weight $2$, and the predicted level is prime to $11$
and divides the conductor of $A$, i.e.\ $\mathfrak p\cdot (2)^{a}$ for some $a\ge0$.

Let $S_2(\mathfrak n)$ be the space of Hilbert cusp forms over $K$ of parallel weight $2$ and
level $\mathfrak n$, computed via the definite quaternion algebra ramified at the two
archimedean places, with its commuting Hecke operators $T_P$ (Brandt matrices), which are
integral. Reduce everything mod $11$. If a Hilbert newform $f$ has
$\rhob_{f,\lambda}\cong\rho_i$ for some degree-one prime $\lambda\mid 11$ of its Hecke field,
then its reduction is an eigenvector $v$ in $S_2(\mathfrak n)\otimes\bbF_{11}$ with
$T_P v = t_P\,v$ where $t_P$ is a root of \eqref{eq:fingerprint}; that is,
$\big(T_P^2 - \Tr_P\,T_P + \Nm_P\big)v=0$.

Empirically the matching forms realize $\rho_i$ only after a {\it quadratic twist}
$\varepsilon$ ramified at $2$, so $t_P$ is replaced by $\varepsilon(P)\,t_P=\pm t_P$. We
therefore intersect the twist-robust kernels: set
\begin{equation}\label{eq:V}
V \;=\; \bigcap_{P}\ \Big[\ \ker\!\big(T_P^2 - \Tr_P\,T_P + \Nm_P\big)\;+\;
\ker\!\big(T_P^2 + \Tr_P\,T_P + \Nm_P\big)\ \Big] \ \subseteq\ S_2(\mathfrak n)\otimes\bbF_{11}.
\end{equation}
A common eigenform realizing $\rho_i$ (up to twist) lies in $V$.

\begin{theorem}[computational]\label{thm:main}
Take $\mathfrak n=\mathfrak p\cdot(2)$, of norm $264716=4\cdot 66179$; then
$\dim S_2(\mathfrak n)=5514$. The space $V$ of \eqref{eq:V}, intersected over the primes $P$
above $\{29,31,41,7,59\}$ (eight prime conditions), is
\[
\dim V \;=\; 2,
\]
stable: $\dim S_2(\mathfrak n)\otimes\bbF_{11}=5514\to 5\to 2\to 2\to\cdots\to 2$. By contrast
$V=0$ at level $\mathfrak p$ alone, and $V=0$ at level $\mathfrak n$ if one keeps only the
untwisted kernel.
\end{theorem}

A $2$-dimensional space surviving eight independent constraints ``$a_P\in\{\pm t_P,\pm t_P'\}$''
is not an accident (each prime cuts the surviving dimension by a definite factor).

Restricting each $T_P$ to $V$ yields a \emph{scalar} matrix --- a single eigenvalue $a_P$ ---
so $V$ is the multiplicity-$2$ eigenspace of \emph{one} mod-$11$ Hecke eigensystem $\sigma$
(two distinct forms would give two eigenvalues, and $\rho_1,\rho_2$ cannot share an eigensystem
since they are not twist-related). Thus the match exhibits \emph{one} of the two pieces,
realized with multiplicity $2$.

Crucially, $\sigma$ \emph{is} a genuine Hecke eigensystem, hence a genuine $2$-dimensional
Galois representation, so it follows one piece consistently: setting
$\mathrm{tr}\,\rho_1(\Frob_P)$ to be the root of $y^2-\Tr_P y+\Nm_P$ whose square equals
$a_P^2$ gives a globally consistent labeling of the two pieces (and $\mathrm{tr}\,\rho_2$ is the
complementary root). Thus the modular form itself separates $\rho_1$ from $\rho_2$ --- no
explicit computation of the polarization kernel $K(\mathcal L)$ is needed.

\begin{theorem}\label{thm:serre}
One of the two $2$-dimensional mod-$11$ representations $\rho_1,\rho_2$ attached to
$\Jac(C)/\bbQ(\sqrt5)$ is modular up to a quadratic twist: the eigensystem $\sigma$ of
Theorem~\ref{thm:main} satisfies $a_P^2=\mathrm{tr}\,\rho_i(\Frob_P)^2$ at all $24$ primes
checked, so $\mathrm{tr}\,\sigma=\varepsilon(P)\,\mathrm{tr}\,\rho_i$ with
$\varepsilon=\chi_{K(\sqrt d)}$, $\ d=\tfrac12\big(330895-66179\sqrt5\big)=66179\cdot
\tfrac{5-\sqrt5}{2}$, a quadratic character ramified at the prime above $5$ and the primes
above $66179$ (unramified at $2$). Since a quadratic twist preserves modularity, this realizes
$\rho_i$ as modular, in agreement with Serre's conjecture over $\bbQ(\sqrt5)$.
\end{theorem}

\begin{remark}
This is computational evidence, not a proof. The dimension count in Theorem~\ref{thm:main} is
unconditional; the eigensystem $\sigma$ is extracted and verified ($a_P^2=\mathrm{tr}\,\rho_i^2$)
over $24$ primes; and $\varepsilon=\chi_{K(\sqrt d)}$ matches $17$ of the $18$ non-degenerate
cleanly-signed primes, the lone exception being one prime above $59$ (a residual extraction
artifact; the excluded $19$th point is the double-root prime above $61$, where $\varepsilon$ is
ill-defined). What remains is to confirm that the \emph{second} piece $\rho_2$ is itself
modular (its traces are now explicit, as the complementary roots) and to remove the one
residual misfit. All constructions are reproducible from the accompanying scripts
(\texttt{mod11\_rep.m}, \texttt{match\_hmf\_twist.m}, \texttt{extract\_eigensystems.m},
\texttt{identify\_twist2.m}, \texttt{verify\_twist.m}; fingerprint in \texttt{mod11\_trnm.txt}).
\end{remark}

\end{document}
